\section{Background and Motivation}\label{sec:intro:background}

Radar and communications systems share a limited electromagnetic spectrum.
Passive \gls{esm} receivers listen without transmitting and record time-ordered \glspl{pdw} that feed \gls{elint} analysis: detecting radars, grouping pulses by emitter, and recognizing operating modes \parencite{wiley2006elint,skolnik2008radar}.

Modern radars vary carrier frequency, pulse width, amplitude, and \gls{pri} from burst to burst.
That agility weakens simple fingerprints and makes classical parsers built for stable pulse trains less reliable.
The same emitter also switches \emph{modes}: recurrent control policies that map a task such as volume search or track to a characteristic pulse regime.
Knowing the mode constrains how the waveform may evolve and can separate emitters that look similar at the type level.

The input to this thesis is a stream of \glspl{pdw}.
Several emitters are interleaved in time, so the encoder must use context across the window.
The supervision available in labelled training data is asymmetric: emitter identifiers are unique only within a recording, while operating modes come from a shared catalogue.

This thesis asks whether a transformer encoder, trained with clustering objectives, can map such streams to embeddings in which emitters and modes both form separable groups.
